\scanpage{96}

D'autre part $x,y$ sont orthogonaux : $x\wedge y$ donc \unsure{à|:} $\deleted{t}\unsure{z|t}=\widetilde f(\unsure{\xi|\zeta})$, qui n'est \deleted{pas} orth. à lui-$=$ en aucune fibre, car $\varphi(t,t)=2q(t)\deleted{\unsure{=}}=-2\ne0$ en chaque fibre. Donc $x,y$ forment une base de $t^\perp$, dito de $=$ $(x',y')$. On doit donc avoir
\[
\left\{\begin{aligned}x'&=\unsure{a}x\unsure{+|-}by\\y'&=cx+dy\end{aligned}\right.\qquad\text{avec }ad-bc\in\Gamma G_m
\]

\deleted{Ce qui donne} Écrivons alors $q(x')=q(y')$, $\varphi(x',y')=0$, on \unsure{trouve}. 
\[
\begin{aligned}\inserted{\varphi(x',y')=}\;ac\deleted{\unsure{+}}\varphi(x,x)-bd\,\varphi(y,y)&=0\quad\text{i.e. comme }\ \begin{aligned}\varphi(x,x)&=2q(x)\\[-.4ex] &\hspace{1.4em}\shortparallel\\[-.4ex]\varphi(y,y)&=2q(y)\end{aligned}\\[-.1ex]\text{inversibles}\qquad ac-bd&=0 \qquad\qquad \begin{matrix}a&b\\d&c\end{matrix}\end{aligned}
\]
donc $\overbrace{c=\lambda b,\ d=\lambda a\deleted{\unsure{,}}}$ avec $\lambda\in\Gamma G_m$

\deleted{et \unsure{on}} \unsure{Or|d'où}
\[
\left\{\begin{aligned}q(x')&=(a^2+b^2)q(x)\\q(y')&=(c^2+d^2)q(y)=\lambda^2(a^2\unsure{+|-}b^2)q(x)\end{aligned}\right.
\]
donc la condition $q(x')=q(y')$ \inserted{$\in\Gamma(G_m)$} signifie :
\[
\deleted{(\unsure{\lambda^2-1|\lambda^2})(a^2+b^2)=0}
\]
\[
\left\{\begin{aligned}a^2+b^2&\in\Gamma G_m\\\lambda^2&=1\end{aligned}\right.
\]

\deleted{donc}\qquad \deleted{$(x',y')=\qquad\unsure{\xi'|\int'}=x'+y'=\lambda$}
% archive sheet number: 95
